% !TeX root = ../informe.tex
% ===========================================================================
%  4. Arquitectura de la solución
%
%  Sección central del informe: la estructura del sistema y, sobre todo, su
%  justificación. El modelo de datos va aparte, en la sección 5.
%
%  Deja las figuras en figuras/ con estos nombres, en PDF:
%      contexto.pdf · contenedores.pdf · componentes.pdf · despliegue.pdf
%  Si alguna falta, el documento compila igual y deja un recuadro visible
%  en su lugar.
% ===========================================================================

\section{Arquitectura de la solución}
\label{sec:arquitectura}

\instruccion{Esta sección responde a una sola pregunta: por qué el sistema
está partido así y no de otra manera. Describir los componentes es la parte
fácil; lo que se evalúa es la justificación. Apóyate en las vistas para
mostrar la estructura y en el registro de decisiones para explicarla.}

\subsection{Visión general}
\label{sec:vision-general}

\ph{Dos o tres párrafos: qué estilo arquitectónico se aplica, en cuántas
piezas se divide el sistema y qué principio gobierna el reparto.}

\subsection{Notación}
\label{sec:notacion}

\instruccion{Di en qué notación están los diagramas y con qué herramienta se
generaron. Si salen de una única descripción textual, dilo: es el argumento
de que las vistas son consistentes entre sí por construcción.}

Los diagramas siguen el modelo C4 \cite{ejemplo-web}: contexto, contenedores
y componentes, más una vista de despliegue. \ph{Herramienta y fuente desde
la que se generan.}

\subsection{Vistas arquitectónicas}
\label{sec:vistas}

\instruccion{Una subsección por vista. En cada una: el diagrama, y debajo la
prosa que explica lo que el diagrama no puede decir — por qué esa frontera,
qué pasa si un elemento falla, qué se decidió a propósito. Un diagrama sin
texto no comunica; un texto que solo enumera lo que ya se ve en el diagrama,
tampoco. Quita las vistas que no apliquen en vez de dejarlas vacías.}

\subsubsection{\texorpdfstring{\textsf{V-01}}{V-01} Contexto del sistema}
\label{v:01}

\begin{figure}[H]
  \centering
  \diagrama{contexto}
  \caption{Vista de contexto: el sistema, sus usuarios y los sistemas
           externos con los que se relaciona.}
  \label{fig:contexto}
\end{figure}

\ph{Qué resuelve el sistema y para quién. Un párrafo.}

\subsubsection{\texorpdfstring{\textsf{V-02}}{V-02} Contenedores}
\label{v:02}

\instruccion{Suele ser la vista principal del documento: es donde se ve la
descomposición completa del sistema en piezas ejecutables. Si el diagrama no
cabe legible en el bloque de texto, \texttt{\textbackslash diagramapagina}
lo lleva a página completa.}

\begin{figure}[p]
  \thisfloatpagestyle{empty}
  \centering
  \diagramapagina{contenedores}
  \caption{Vista de contenedores: piezas ejecutables del sistema y sus
           almacenes de datos.}
  \label{fig:contenedores}
\end{figure}

\begin{table}[H]
  \centering
  \caption{Composición del sistema.}
  \label{tab:componentes}
  \begin{tabular}{L{3.4cm} L{6.4cm} L{3.6cm}}
    \toprule
    \textbf{Pieza} & \textbf{Responsabilidad} & \textbf{Tecnología} \\
    \midrule
    \id{\ph{...}} & \ph{...} & \ph{...} \\
    \id{\ph{...}} & \ph{...} & \ph{...} \\
    \id{\ph{...}} & \ph{...} & \ph{...} \\
    \bottomrule
  \end{tabular}
\end{table}

\ph{Recorrido de la tabla: qué agrupa cada pieza y por qué la frontera cae
donde cae. Si alguna no tiene almacén propio, explica de dónde saca los
datos.}

\paragraph*{Comunicación.}
\ph{Qué pieza llama a cuál, con qué protocolo, y qué ocurre si el destino no
responde.}

\subsubsection{\texorpdfstring{\textsf{V-03}}{V-03} Componentes}
\label{v:03}

\instruccion{Detalla una sola pieza, la que concentre la lógica de negocio,
y di explícitamente por qué esa y no otra. Un diagrama de componentes por
cada pieza infla el documento sin añadir nada.}

\begin{figure}[p]
  \thisfloatpagestyle{empty}
  \centering
  \diagramapagina{componentes}
  \caption{Vista de componentes de \id{\ph{pieza}}.}
  \label{fig:componentes}
\end{figure}

\ph{Recorrido de una petición representativa a través de los componentes.}

\subsubsection{\texorpdfstring{\textsf{V-04}}{V-04} Despliegue}
\label{v:04}

\begin{figure}[p]
  \thisfloatpagestyle{empty}
  \centering
  \diagramapagina{despliegue}
  \caption{Vista de despliegue: nodos de ejecución, redes y volúmenes.}
  \label{fig:despliegue}
\end{figure}

\ph{Cómo se materializa la arquitectura en tiempo de ejecución: qué nodos y
redes hay, qué aísla cada una y dónde persisten los datos. El detalle
operativo (comandos, variables de entorno) va en el manual de despliegue, no
aquí: duplicarlo garantiza que los dos se desincronicen.}

\subsection{Seguridad y control de acceso}
\label{sec:seguridad}

\ph{Dónde se autentica el usuario, quién emite la credencial, quién la
valida y cómo llega la identidad a cada pieza del sistema.}

\subsection{Registro de decisiones de arquitectura}
\label{sec:decisiones}

\instruccion{Cada decisión relevante con su alternativa descartada. Una
decisión sin alternativa no es una decisión: es una descripción. Este
registro es lo que distingue un documento de arquitectura de un manual de
usuario del propio sistema. Referencia las decisiones desde el texto con
\texttt{\textbackslash adrref\{01\}}. Repite el bloque una vez por decisión.}

\decision{01}{\ph{Título de la decisión}}
  {\ph{Qué problema o fuerza obligó a decidir.}}
  {\ph{Qué se decidió.}}
  {\ph{Qué otra opción se consideró y por qué se descartó.}}
  {\ph{Qué se gana y qué se paga por ello.}}

\instruccion{Si no sabes por dónde empezar, las decisiones que casi siempre
hay que justificar son: cómo se reparte la persistencia entre piezas; dónde
se validan las reglas de negocio; si el acceso pasa por un punto único o
cada cliente habla directamente con cada servicio; y qué se resuelve en
tiempo de construcción frente a tiempo de ejecución.}
